%*****************************************
% Set tocdepth to -1 (part level only) so chapters are suppressed
\newcounter{saveddepth}
% Save current depth and write the lowered value to .toc
\setcounter{saveddepth}{\value{tocdepth}}
\addtocontents{toc}{\protect\setcounter{tocdepth}{-1}}
\addtocounter{chapter}{-1}
\renewcommand{\theHchapter}{\thechapter.repeated} % Provides a unique internal anchor
\chapter{Conclusiones}\label{ch:conclusiones}
%*****************************************

\begin{flushquote}
    {Essentially, all models are wrong, but some are useful.}
    {George E. P. Box}
\end{flushquote}

A lo largo de este trabajo hemos desarrollado un marco bayesiano para inferencia con datos funcionales basado en \acs{RKHS}. La idea principal consiste en aprovechar la estrecha relación entre un proceso estocástico y el \acs{RKHS} generado por su función de covarianza, con el fin de construir modelos funcionales que sean a la vez matemáticamente robustos y computacionalmente viables. Al representar el parámetro funcional como una combinación lineal finita de secciones del núcleo, obtenemos modelos que conservan una interpretación funcional, mientras que reducimos el problema de realizar inferencia bayesiana en un espacio de dimensión infinita a un problema paramétrico de dimensión variable.

Dentro de este marco hemos abordado dos problemas estadísticos fundamentales. El primero es el de regresión funcional, para el cual se han propuesto modelos lineales y logísticos bayesianos, junto con estrategias de predicción y procedimientos de selección de variables. El segundo es el contraste de hipótesis funcionales, donde se ha construido un test bayesiano de una muestra para la media de un proceso gaussiano mediante una distribución a priori \textit{spike-and-slab}. Esta metodología permite tomar una decisión respecto a la hipótesis nula, pero también proporciona medidas a posteriori de la evidencia en su contra y una cuantificación directa de la incertidumbre asociada al efecto funcional subyacente.

La faceta bayesiana de nuestros modelos cumple varias funciones. En primer lugar, proporciona un procedimiento de inferencia sencillo y flexible para modelos basados en \acs{RKHS}, aplicable a una amplia gama de distribuciones a priori y, en general, interesante desde el punto de vista práctico. En segundo lugar, mejora la interpretabilidad de las decisiones obtenidas, ya que da acceso a una distribución a posteriori completa que describe la influencia de cada parámetro, así como las interacciones entre ellos. En tercer lugar, permite cuantificar de forma natural la incertidumbre del proceso, aportando información adicional mediante resúmenes a posteriori o bandas creíbles.

Una parte importante de la tesis ha consistido en establecer los fundamentos teóricos de estos procedimientos bayesianos. Para la regresión funcional, hemos obtenido consistencia a posteriori en presencia de un número desconocido de componentes basándonos en el teorema de Doob, junto con resultados sobre tasas de convergencia siguiendo el enfoque más sofisticado de Schwartz. Otros autores, como \citet{lian2016posterior} y \citet{abraham2020posterior}, han desarrollado análisis similares, aunque no dentro de un marco basado en \acs{RKHS}. En el caso de los contrastes funcionales, el trabajo de \citet{dass2004note} estableció la consistencia del factor de Bayes en problemas de contraste con una hipótesis nula puntual, pero su metodología abstracta no permite obtener tasas de contracción a posteriori. Gracias al enfoque \acs{RKHS} hemos podido obtener resultados más completos que incluyen estas tasas de convergencia, caracterizando además la diferencia entre la tasa polinómica bajo la hipótesis nula y la tasa exponencial bajo la hipótesis alternativa. También hemos establecido que la tasa exponencial obtenida se encuentra arbitrariamente cerca de la tasa óptima. Estos resultados proporcionan una justificación frecuentista para nuestro test bayesiano, a la vez que aclaran el comportamiento asintótico de la evidencia acumulada por los modelos.

Asimismo, el carácter transdimensional de la aproximación de la distribución a posteriori constituye una contribución novedosa en el ámbito de los datos funcionales. Al combinar la representación \acs{RKHS} con el muestreo \acs{RJMCMC}, la complejidad de los modelos puede inferirse conjuntamente con sus parámetros, sin necesidad de una etapa independiente de selección de modelos. Este enfoque sigue la tendencia reciente en metodología bayesiana, donde una mayor capacidad computacional permite utilizar en la práctica técnicas de aproximación más costosas. El rendimiento obtenido en los experimentos numéricos es generalmente satisfactorio, incluso con el número relativamente pequeño de componentes que hemos seleccionado. En el contexto de la regresión, los métodos basados en \acs{RKHS} ofrecen un rendimiento predictivo comparable al de un amplio conjunto de alternativas funcionales y multivariantes. En el contexto de los contrastes de hipótesis, nuestros métodos proporcionan un control preciso del error de tipo I y una potencia adecuada en diversas situaciones.

En conjunto, hemos desarrollado un enfoque bayesiano unificado, interpretable y sustentado teóricamente para inferencia funcional, que combina representaciones mediante \acs{RKHS}, aproximación transdimensional de la posteriori, cuantificación de la incertidumbre y garantías asintóticas rigurosas. Los modelos y métodos propuestos no pretenden sustituir a las técnicas existentes, sino complementarlas mediante alternativas que ofrecen un equilibrio diferente entre flexibilidad, interpretabilidad, viabilidad computacional y fundamento teórico.

\section{Trabajo futuro y líneas de investigación abiertas}

Consideramos que este marco bayesiano abre diversas líneas de investigación prometedoras y pone de manifiesto tanto la utilidad práctica como el interés teórico de las representaciones mediante \acs{RKHS} para abordar la naturaleza intrínsecamente infinito-dimensional de los datos funcionales. A continuación, presentamos algunos ejemplos de posibles direcciones para futuros trabajos.

\begin{itemize}
    \item El estudio del papel de los puntos de impacto dentro de la aproximación mediante \acs{RKHS} podría dar lugar a nuevas metodologías para la detección de puntos de cambio o la selección de variables en estimación funcional. Por ejemplo, \citet{li2021bayesian} descomponen las observaciones funcionales en coeficientes \textit{wavelet} e introducen un parámetro de punto de cambio específico para cada coeficiente, asignando distribuciones a priori independientes a los valores de los coeficientes antes y después del cambio. Una posibilidad sería considerar un modelo alternativo en el que la descomposición se sustituyera por una aproximación basada en \acs{RKHS}.
    \item El análisis teórico de los modelos de regresión funcional todavía requiere resultados explícitos para el caso logístico. Cabe esperar que sean similares a los obtenidos para el caso lineal, aunque las demostraciones podrían ser más complejas y requerir técnicas diferentes que permitan comprender mejor las particularidades de la formulación basada en \acs{RKHS}.
    \item La metodología de contraste de hipótesis podría extenderse a datos no gaussianos mediante mixturas de gaussianas o mediante una distribución más general. Una extensión factible consistiría en introducir etiquetas latentes \(Z_i\in\{1,\dots,M\}\) y suponer \(\pp(Z_i=m)=w_m\) con \(\sum_m w_m=1\), para posteriormente considerar, en el caso de una covarianza común, el modelo \(X_i| Z_i=m\sim\gp(\mu_m,K)\). Cada función de media \(\mu_m\) se aproximaría entonces mediante una representación \acs{RKHS} de dimensión finita.
    \item Nuestro test de una muestra cubre de forma natural el caso de datos funcionales emparejados, que es bastante habitual, por ejemplo, en estudios sobre efectos de tratamientos o fármacos. Una extensión natural consistiría en considerar el caso general de contrastes de dos muestras, evaluando la igualdad de las funciones medias de dos poblaciones funcionales independientes, tanto en el caso homocedástico como en el heterocedástico.
    \item Explorar distintas elecciones para la distribución a priori de los coeficientes en el contraste funcional podría mejorar las tasas de convergencia. Por ejemplo, \citet{johnson2010use} proponen distribuciones a priori <<no locales>> que excluyen explícitamente aquellas densidades a priori alternativas que son positivas en \(\Theta_0\). Otro ejemplo es el trabajo de \citet{rockova2018bayesian}, que considera una extensión continua de la distribución \textit{spike-and-slab} en la que el \textit{spike} se concentra en torno a \(b=0\), pero no es una delta de Dirac.
    \item Desde el punto de vista computacional, existe un amplio margen para mejorar los mecanismos de propuesta de nuevos parámetros en la etapa de incremento de dimensión del algoritmo \acs{RJMCMC}.
\end{itemize}
%*****************************************
%*****************************************
%*****************************************
%*****************************************
%*****************************************
% Restore tocdepth
\addtocontents{toc}{\protect\setcounter{tocdepth}{\thesaveddepth}}